% This section measures the precision of long-form generations from the state-of-the-art LMs using human evaluation.

We introduce \ours, a new evaluation of an LM that considers the factual precision of atomic facts generated by the LM.
We perform human evaluations to calculate \ourScores\ of the state-of-the-art LMs (Section~\ref{subsec:data}) and discuss results (Section~\ref{subsec:annotation-results}).
%
\ours\ allows rigorous and fine-grained evaluation of factual precision, but is time-consuming and costly, motivating automatic evaluation in Section~\ref{sec:method}.


\subsection{Definition}
\label{subsec:data-overview}

\ours\ is based on two key ideas.

\vspace{.3em}
\noindent
\textbf{Key idea 1: Atomic fact as a unit.}
Long-form text consists of many pieces of information that can each be either true or false.
Prior work has explored using a sentence as a unit; however, even a single sentence is a mix of supported and unsupported facts, e.g., in 40\% of the cases with ChatGPT.
Previous and concurrent work either (1) defines an additional label of \texttt{partial support}~\citep{manakul2023selfcheckgpt,liu2023evaluating} whose definition may be subjective and can lead to low agreement, %is inherently not well-defined 
%\pw{``inherently'' is doing heavy lifting here; why is this not inherently well-defined?}
or (2) takes the strictest definition of \texttt{support} that requires every piece of information to be supported~\citep{rashkin2021measuring,gao2022attributed}, which ignores the partial support cases, e.g., assigning 0.0 to both generations in Figure~\ref{fig:teaser} even though the first generation is considerably more accurate than the second.
%in an example in Figure~\ref{fig:teaser}, this approach assigns 0.0 to both generations even though the first generation is considerably more accurate than the second.
\myskip{
\begin{table}[t]
    \footnotesize \setlength{\tabcolsep}{0em}
    \centering
    \begin{tabular}{p{1.2cm}p{6cm}}
        \toprule
        \multicolumn{2}{l}{
        \textbf{Sentence:} Thierry Henry (born 17 August 1977) is a French} \\
        \multicolumn{2}{l}{professional football coach, pundit, and former player.} \\
        \midrule
        \textbf{Fact 1:} &    
        Thierry Henry was born on 17 August 1977. \\
        \textbf{Fact 2:} &    
        Thierry Henry is French. \\
        \textbf{Fact 3:} &
        Thierry Henry is a professional football coach. \\
        \textbf{Fact 4:} &
        Thierry Henry is a football pundit.  \\
        \textbf{Fact 5:} &
        Thierry Henry is a former football player. \\
        \bottomrule
    \end{tabular}
    \caption{
      Examples of our process of decomposing sentences into \emph{atomic facts}, each conveying a single piece of information.
      % \pw{How is this different from the two examples of decomposition shown in Fig 1? I wasn't sure what new information it was conveying.} \sewon{mainly because Fig 1 is far and the text is small, and having more examples is better? I agree we can cut it for the 8p version...}
    }\label{tab:example-atomic-fact-decomposition}
\end{table}}

In this paper, we define an atomic fact as a short sentence conveying one piece of information (examples in Figure~\ref{fig:teaser}), similar to summarization content units~\citep{nenkova2004evaluating}.
%
An atomic fact is a more fundamental unit than a sentence for a piece of information and provides a more fine-grained evaluation, e.g., in Figure~\ref{fig:teaser}, rating the first generation higher than the second.

\vspace{.3em}
\noindent
\textbf{Key Idea 2: Factual precision as a function of a given knowledge source.}
Prior work often considers factual precision as a single global truth~\citep{manakul2023selfcheckgpt}.
In contrast, we adopt a perspective that the truthfulness of a statement should depend on a particular knowledge source that end users consider to be trustworthy and reliable.
%In contrast, we take a perspective that what is counted as {\em true} or {\em false} is inherently ill-defined \pw{(is it inherently ill-defined? or is just that some sources conflict because they're not reliable? it seems like factual statements like ``X was born in Y'' have a well-defined truth value, and we say later that ``Human biographies are objective (not subjective or debatable)''} \sewon{I think my perspective is that, we know ``X was born in Y'' because some evidence shows it. We don't actually know if it is true.} \kalpesh{How about we drop ill-defined and reword it to "In contrast, we adopt a perspective that the truthfulness of every statement should depend on a particular knowledge source that end users consider to be trustworthy and reliable."}
%and should generally depend on knowledge sources that end users consider to be trustworthy and reliable.
Therefore, instead of whether an atomic fact is globally true or false, we consider whether it is {\em supported} by a given source of knowledge.
This has been used in the fact verification literature~\citep{wadden2022scifact} where conflict of information between different sources is relatively common.
%A similar definition has been previously adopted for evaluating faithfulness in summarization, although at a sentence level~\citep{pagnoni-etal-2021-understanding,pasunuru-etal-2021-efficiently,huang-etal-2021-efficient}.

\vspace{.3em}
\noindent
\textbf{Definition.}
Let $\mathcal{M}$ be a language model to be evaluated,
$\mathcal{X}$ be a set of prompts,
and $\mathcal{C}$ be a knowledge source.
Consider a response $y = \mathcal{M}_x$ for $x \in \mathcal{X}$
and $\mathcal{A}_y$, a list of atomic facts in $y$.
A \ours\ of $\mathcal{M}$ is defined as follows.
\begin{align*}
    \begin{gathered}
        f(y) = \frac{1}{|\mathcal{A}_y|}\sum_{a \in \mathcal{A}_y}\mathbb{I}[a \text{~is supported by~} \mathcal{C}], \\
        \text{\ours}(\mathcal{M}) =
        \mathbb{E}_{x \in \mathcal{X}}[f(\mathcal{M}_x) | \mathcal{M}_x\text{~responds}].
    \end{gathered}
\end{align*} {\em $\mathcal{M}_x$ responds} means $\mathcal{M}$ did not abstain from responding to the prompt $x$.
% \pw{$p$ is not meant to denote probability, right? perhaps we can pick a different letter to avoid confusion?}
\myskip{
\vspace{10em}
Let $x$ be a long-form generation from an LM that conveys a series of atomic facts, i.e., pieces of information, $x_1, x_2, \cdots, x_n$, and $\mathcal{C}$ be a knowledge source.
%We define \textbf{factual precision} $f$ as a function that maps $x$ and $\mathcal{C}$ into:
We define \textbf{\ours} as a function that maps $x$ and $\mathcal{C}$ into:
\begin{equation*}
    \text{\ours}(x, \mathcal{C}) = \frac{1}{n} \sum_{i=1}^n \mathbb{I}[x_i \text{~is supported by~} \mathcal{C}].   
\end{equation*}
\begin{equation*}
    \text{\ours}(\mathcal{M}, \mathcal{C}) = \frac{
        \sum_{x \in \mathcal{X}} f(x, \mathcal{C}) \mathrm{DoesRespond}(x)
    }{
        \sum_{x \in \mathcal{X}} \mathrm{DoesRespond}(x)
    }
\end{equation*}}
% A similar definition has been used in evaluating faithfulness in summarization~\citep{pagnoni-etal-2021-understanding,pasunuru-etal-2021-efficiently,huang-etal-2021-efficient}.
This definition assumes the following:
\vspace{-.3em}
\begin{enumerate}[leftmargin=15pt]\itemsep -.2em
    \item Whether or not an atomic fact is supported by $\mathcal{C}$ is undebatable.
    \item Every atomic fact in $A_y$ has an equal weight of importance, following \citet{krishna-etal-2023-longeval}. 
    \item Pieces of information in $\mathcal{C}$ do not conflict or overlap with each other.
\end{enumerate}
In the rest of the paper, we propose to use people biographies as $\mathcal{X}$ and Wikipedia as $\mathcal{C}$ because they satisfy these assumptions to a reasonable degree (Section~\ref{subsec:data}).
We discuss in which cases these assumptions hold or may not hold in more detail in the Limitation section. % \pw{forward ref} \sewon{We don't have a section number for the limitation section... :joy:}

\ours\ considers {\em precision} but not {\em recall}, e.g., a model that abstains from answering too often or generates text with fewer facts may have a higher \ours, even if these are not desired.
We leave the evaluation of factual recall for future work (more discussion in the Limitation section). % and recommend using \ours\ in conjunction with the average number of atomic facts produced.
%\mohit{should you mention somewhere that LMs that generate more generic text (fewer specific details) are likely to have inflated factscores? like in fig 1, what if an LM just generated ``Bridget Moynahan is an American actress''. this is different than abstaining, and it would be rewarded with a perfect factscore despite not providing more specifics on what movies/shows she acted in etc. if factscore becomes a metric to eval/develop LMs, it could end up preferring shorter / more generic generations.} \sewon{great point. wondering if the above text addresses this?}
% \sewon{Previous version here:}


% We break a generation from an LM into a series of atomic facts,
% i.e., short sentences, each conveying one piece of information, similar to summarization content units (SCUs) in the summarization literature~\citep{nenkova2004evaluating}. We then assign a precision label to each atomic fact. This is in contrast to recent work that assigns a sentence-level label~\citep{gao2022attributed,manakul2023selfcheckgpt,liu2023evaluating}.
% Assigning finer-grained labels is important for a few reasons.\kalpesh{below will be stronger with more supporting citations beyond SCUs.. FRANK and LongEval are somewhat relevant, but atomicity of facts was not the primary focus in those works.. I'm wondering if there's some HCI literature on this? HCI citations will definitely be compelling here}
% \vspace{-.3em}
% \begin{enumerate}[leftmargin=15pt]\itemsep -.2em
%     \item Precision labels are better defined. Each sentence contains multiple pieces of information, often a mix of supported and unsupported facts. Prior work dealt with this issue by assigning \texttt{support} or \texttt{partial support}, which does not fundamentally address the problem and leads to a higher disagreement rate. 
%     \item It leads to better annotation quality. Fact-checkers often miss minor or implicit facts when making a sentence-level judgment, but judge more comprehensively when they are given a set of atomic facts to be validated.
%     \item When the goal is to make a sentence-level judgment, validating each piece of information first and aggregating the results is still a more natural and better way of making a sentence-level judgment, for both humans and models. In Section~\ref{subsec:models-validation}, we show that our validation models that follow such a pipeline are significantly better than prior approaches that directly make a sentence-level judgment.
% \end{enumerate}
% \subsection{\hanna{title: Studied LMs:}\subjectLM{}}
\subsection{Studied LMs}
We evaluate three LMs (referred to as \subjectLM{}, %\footnote{To distinguish from an LM as an evaluator (Section~\ref{sec:method}).} 
an LM as a subject):
%We consider three LMs that are being evaluated 
%\pw{We evaluate three LMs:}
%(referred to as \subjectLM{}\footnote{To distinguish from an LM as an evaluator (Section~\ref{sec:method}). \pw{could we use the full word that we're abbreviating somewhere in the sentence? Or spell it out in the subsection title?}}):
(1) \textbf{InstructGPT} (\texttt{text-davinci-003}, updated from \citet{ouyang2022training}),
(2) \textbf{ChatGPT}~\citep{chatgpt}, %the most recent model from OpenAI at the time of conducting experiments,
and (3) \textbf{PerplexityAI},\footnoteref{fn:ppl} which incorporates a search engine with a language model. %\footnote{\citet{clarityai} estimates that it feeds snippets from a commercial search engine to the OpenAI API with instructions to only generate information provided in search results, along with citations. \citet{liu2023evaluating} reports that PerplexityAI provides the top utilities and citations over a range of other search-incorporated LMs.}

\subsection{Data}
\label{subsec:data}


We perform human evaluation of factual precision based on our definition. We prompt the \subjectLM{} to generate {\em people biographies} and evaluate them against Wikipedia for the following reasons.
\vspace{-.3em}
\begin{itemize}[leftmargin=15pt]\itemsep -.2em
    \item Biographies are objective (not subjective or debatable) and contain specific (not vague) information, satisfying Assumption 1 in Section~\ref{subsec:data-overview}. %\footnote{Although {\em model-generated} biographies are not guaranteed to be so, we find it is generally the case with (near) state-of-the-art LMs, likely because LMs have been exposed to many human biographies during pre-training.}
    \item Biographies allow evaluation across diverse nationalities, professions, and levels of rarities.
    \item Wikipedia offers reasonable coverage of information about people and is reasonably self-consistent,\footnote{See Appendix~\ref{subsec:annotation-qualitative} for a related analysis.} satisfying Assumption 3. % in Section~\ref{subsec:data-overview}.
\end{itemize}
%We additionally note that the generation of a people bio is a fairly easy task, and simply copying the first paragraph of a Wikipedia page about a person guarantees perfect factual precision.


\paragraph{Data collection.}
We carefully design an annotation pipeline to assign a factual precision to a long-form generation through the following steps.

\vspace{.4em}
\noindent \textbf{Step 0: Sampling people entities.}
We sample 183 people entities from Wikidata who have %the \pw{-the}
corresponding Wikipedia pages.
We sample entities to annotate from a uniform distribution over categories defined in Appendix~\ref{app:data-category}. %\footnote{A concurrent work~\citep{manakul2023selfcheckgpt} also collected the biography data focusing on the frequent entities whose Wikipedia pages are long.  This difference leads to significant differences in factual precision as discussed in Section~\ref{subsec:annotation-results}.}

\vspace{.4em}
\noindent \textbf{Step 1: Obtaining generations.} We feed a prompt \texttt{``Tell me a bio of <entity>''} to the \subjectLM{} and take a generation as it is. %\footnote{While different instructions or sampling may lead to different results, we assume every \subjectLM\ is a black box and take a generation with a minimal prompt. In fact, LMs are often being served with additional instructions hidden to end users to provide the best possible quality, e.g., PerplexityAI, chat-based LMs in Section~\ref{sec:eval-new-lms}.}
We implement rules to identify generations that abstain from answering and filter them out. % and send the rest of the generations to the next steps.

\vspace{.4em}
\noindent \textbf{Step 2: Atomic facts generation.}
Human annotators break a generation into a series of atomic facts.
To save annotation time, we provide atomic facts broken down by InstructGPT which human annotators can take and revise. Details in Appendix~\ref{app:atomic-fact-generation}.

\vspace{.4em}
\noindent \textbf{Step 3: Labeling factual precision \& editing.}
We ask another set of human annotators to assign each atomic fact one of three labels. If the atomic fact is clearly not related to the prompt, and thus should be removed from the bio without a validation step, they assign \irLabel. If the fact is relevant, they validate the fact based on the English Wikipedia, and label either \sLabel\ or \nsLabel. % Additionally, annotators also edit the text to make it factually correct. They are asked to correct factual errors or remove the sentence if the information in the sentence is entirely off. They are asked not to revise the text to improve the coverage of the information or dimensions beyond factual precision. \pw{what do we use this for?}

\vspace{.3em}
We recruit freelancers through Upwork and pay 15--25 USD per hour. Annotation requires extensive effort and time, leading to the cost of \$4 per generation. %e.g., completing Step 3 for three LMs takes 25 minutes on average.
We assign two freelancers for the 10\% of the data and calculate the agreement rate: 96\%, 90\% and 88\% for InstructGPT, ChatGPT and PerplexityAI, respectively. More details are provided in Appendix~\ref{app:annotator-recruitment}.

\begin{table}[t]
    \footnotesize \centering
    \begin{tabular}{l @{\hspace{0.5em}} r @{\hspace{1em}} r @{\hspace{1em}} r}
        \toprule
            & InstGPT & ChatGPT & PPLAI \\
        \midrule
            Use search & \xmark & \xmark & \cmark \\
            \% responding & 99.5 & 85.8 & 90.7 \\
            \# tokens / response & 110.6 & 154.5 & 151.0 \\
            \# sentences / response & 6.2 & 7.9 & 9.8 \\
            \# facts / response & 26.3 & 34.7 & 40.8 \\
        \midrule
            \multicolumn{4}{c}{\em Statistics of the labels} \\
                \sLabel & 42.3 & 50.0 & 64.9 \\
                \nsLabel           & 43.2 & 27.5 & 11.1  \\
            \irLabel              
                & 14.0 & 8.3 & 14.8 \\
            Abstains from answering & 0.5 & 14.2  & 9.3 \\
        \midrule
            \textbf{\ours} & \textbf{42.5} & \textbf{58.3} & \textbf{71.5} \\
        \bottomrule
    \end{tabular}
    \caption{
      Statistics of the data and \ours{} results.
      InstGPT and PPLAI respectively refer to InstructGPT and PerplexityAI.
      {\em \% responding} indicates \% of generations that do not abstain from responding.
      {\em \# tokens} is based on white space. % tokenization.
      % Statistics of the labels are macro-averaged over prompts.
      % \textbf{Bold numbers} indicate \ourScores\ defined in Section~\ref{subsec:data-overview}.
    }\label{tab:longform-data-statistics}
\end{table}
\myskip{
\begin{figure}[t]
\centering \footnotesize
\resizebox{1.7\columnwidth}{!}{
    \includegraphics[trim={5.3cm 8cm 6cm 1.7cm},clip]{figures/precision_frequency.pdf}%
}
\resizebox{1.7\columnwidth}{!}{
    \includegraphics[trim={5.3cm 8cm 6cm 2.7cm},clip]{figures/precision_position.pdf}%
}\vspace{-1.5em}
\caption{
    \textbf{Top:} \% of {\em Supported} across varying frequency levels of human entities.
    \textbf{Bottom:} \% of {\em Supported} over relative positions in generation.
    There are significantly fewer supported facts (more precision errors) as the rarity of the entities increases and the position of the fact is later.
}\label{fig:benchmark-ablation}
\end{figure}
}

\subsection{Results}\label{subsec:annotation-results}
Statistics of the data and results are reported in Table~\ref{tab:longform-data-statistics}.

\vspace{-.3em}
\paragraph{All \subjectLM{}s struggle with factual precision errors.}
InstructGPT and ChatGPT achieve \ours{}s of 42.5\% and 58.3\%, respectively.
%InstructGPT and ChatGPT generate supported facts for 42.7--50.5\% of the cases, and supported sentences for 20.4--28.3\%.
PerplexityAI, which uses a commercial search engine and thus should have a perfect \ours\ if directly copying %\pw{grammar}
the text from the correct Wikipedia page, attains a \ours\ of 71.5\%.
We provide a qualitative analysis of its error cases in the last paragraph of this section.
% We qualitatively analyze their error cases in the last paragraph of this section. %Appendix~\ref{subsec:annotation-qualitative}.

ChatGPT and PerplexityAI often abstain from answering which presumably improves their factual precision.
InstructGPT rarely abstains from answering, likely because it is not trained to do so.

Irrelevant facts either
(a) have dependencies on previous facts in a generation that turn out to be unsupported, 
or (b) are irrelevant to the prompt independent from other facts in a generation (examples in Appendix~\ref{app:example-annotated-data}).
We find that (b) rarely happens with InstructGPT and ChatGPT but happens considerably with PerplexityAI, because PerplexityAI often directly copies search results even if they are largely irrelevant to the input prompt.
This is in agreement with a concurrent work from \citet{liu2023evaluating} that shows generative search engines like PerplexityAI copy incorrect search results and generate text that is irrelevant to the input query.

% Irrelevant facts are significantly more frequent in PerplexityAI than in the other two LMs, because PerplexityAI often directly copies search results even if they are largely irrelevant to the input prompt. This is in agreement with a concurrent work from \citet{liu2023evaluating} that shows generative search engines like PerplexityAI copy incorrect search results and generate text that is irrelevant to the input query.

\myskip{
\vspace{-.3em}
\paragraph{Irrelevant facts are frequent in PerplexityAI.}
Irrelevant facts are those that either
(a) have dependencies to previous facts in a generation that turn out to be unsupported, %expand previous facts that turn out to be unsupported,
or (b) are irrelevant to the prompt independent from other facts in a generation (examples in Appendix~\ref{app:example-annotated-data}).
We find that (b) rarely happens with InstructGPT and ChatGPT, but happens considerably with PerplexityAI, because PerplexityAI often directly copies search results even if they are largely irrelevant to the input prompt.
This is in agreement with a concurrent work from \citet{liu2023evaluating} that shows generative search engines like PerplexityAI copy incorrect search results and generate text that is irrelevant to the input query.
}

\begin{table*}[t]
    \center  \myfontsize %\footnotesize
    \setlength{\tabcolsep}{0.5em}
    \begin{tabular}{p{2.3cm}p{0.4cm}p{12.5cm}}
        \toprule
            Category & \% & Example \\
        \midrule
            Single-sentence contradiction & 33.3 &
            \myblue{Gen} On November 25th, 2023, Glover Teixeira became an American citizen. \mypink{Wiki} In November 2020, Teixeira became an American citizen.\\
            (words) && \myblue{Gen} [Eric Hacker] was named the International League Pitcher of the Year. \mypink{Wiki} [Eric Hacker] was named the IL Pitcher of the Week. \\
        \cmidrule(lr){1-3}
            Single-sentence contradiction  & 10.0 &
            \myblue{Gen} William Waldegrave's grandfather was James II and VII. \mypink{Wiki} His father's title was created ... for the diplomat and ambassador James Waldegrave, 1st Earl Waldegrave, whose grandfather was James II and VII.
            \\
            (beyond words) && \myblue{Gen} She has appeared in several successful films such as (...) and Zero (2018). \mypink{Wiki}: Zero was a commercial failure. \\
        \cmidrule(lr){1-3}
            Page-level contradiction & 23.3 & \myblue{Gen} Some of [Julia Faye's] notable films include ... "Cleopatra" (1934). \mypurple{Comment} No mention of {\em Cleopatra} on the {\em Julia Faye} page, and no mention of {\em Julia Faye} on the {\em Cleopatra} page. \\
            && \myblue{Gen} [Kang Ji-hwan] has donated money to various charities and organizations over the years. \mypurple{Comment} No such mention on the {\em Kang Ji-hwan} page.
            \\
        \cmidrule(lr){1-3}
            Subjective & 16.7 & \myblue{Gen} His achievements, as an actor and as a cultural force, will surely prove to be as heroic as those of the characters he portrayed. \mypink{Wiki} Culture writer Steve Rose, in The Guardian, wrote:
            ``Chadwick Boseman began his career playing African American icons and pioneers; he ends it as one himself. His [...] achievements, as an actor and as a cultural force, will surely prove to be as heroic as those of the characters he portrayed.''\\
        \cmidrule(lr){1-3}
            Fact is irrelevant & 3.3 & \myblue{Gen} [Zamfir Arbore]'s life is not well-documented, and there is little information available about him.\\
        \cmidrule(lr){1-3}
            Wiki is inconsistent \& wrong & 3.3 & \myblue{Gen} Kick (2014) that brought [Sajid Nadiadwala] various debutant director awards. \mypink{Wiki} 2015, IIFA Award for Debut Director, Kick. (...) Kick brought him various debutant director awards. \mypurple{Comment} The first text is from a table that indicates he won one award (accurate). The second is inaccurate, incorrectly citing a news article. \\
        \cmidrule(lr){1-3}
            Annotation error & 10.0 &  \myblue{Gen} [Zamfir Arbore] was part of the staff of Românul. \mypink{Wiki} The Românul staff came to include Zamfir Arbore. \mypurple{Comment} Mentioned in the {\em Românul} page but not in the {\em Zamfir Arbore} page. \\
        \bottomrule
    \end{tabular}
    \caption{Categorization of precision errors (\nsLabel) from PerplexityAI (Section~\ref{subsec:annotation-qualitative}).
    \myblue{Gen} indicates the generation from PerplexityAI, and \mypink{Wiki} indicates evidence text from Wikipedia. \mypurple{Comment} indicates our comments.
    }\label{tab:ppl-error-analysis}
\end{table*}

\begin{figure}[t]
\centering \footnotesize
\resizebox{\columnwidth}{!}{
    \includegraphics[trim={22.5cm 8cm 6cm 1.7cm},clip]{figures/precision_frequency_small.pdf}%
}
\resizebox{\columnwidth}{!}{
    \includegraphics[trim={22.5cm 8cm 6cm 2.7cm},clip]{figures/precision_position_small.pdf}%
}\vspace{-1.5em}
\caption{
    \ours\ across varying frequency levels of human entities (\textbf{top}) and relative positions in a generation (\textbf{bottom}). \ourScores\ are lower as the rarity of the entities increases and the position of the fact is later.
}\label{fig:benchmark-ablation}
\end{figure}

\vspace{-.3em}
\paragraph{Error rates are higher for rarer entities.} 
Figure~\ref{fig:benchmark-ablation} (top) shows factual precision over varying frequency levels of topic entities (humans) in the pretraining corpora (see Appendix~\ref{app:data-category}).
There is a notable decrease in \ours\ as the rarity of entities increases, consistently across all \subjectLM{}s.
%There is a significant drop as the entities get rarer.
This is in agreement with \citet{kandpal2022large} and \citet{mallen2022not} which show that short question answering (QA) accuracy is highly correlated with the entity frequencies in the pretraining data. %Drop in the long-form generation is overall more significant than in QA (e.g., an error rate %\pw{an error rate?} increases by 4.0x, compared to 1.8x in \citealp{kandpal2022large}).
However, in contrast to \citet{kandpal2022large} and \citet{mallen2022not} who report QA accuracy of models with retrieval 
is robust to the rarity of entities,
\ours\ of PerplexityAI still significantly drops as entities are rarer: a relative drop of 50\% and 64\% observed at the atomic-level and sentence-level, respectively.

\vspace{-.3em}
\paragraph{Error rates are higher for facts mentioned later in the generation.} 
Figure~\ref{fig:benchmark-ablation} (bottom) reports factual precision over relative positions in a generation. %, e.g., 0--20\% indicates the first 20\% part of the generation.
Across all LMs, the later part of the generation has significantly worse precision. This is likely because (a) information mentioned earlier is more frequently mentioned in the pretraining data (e.g., nationality, profession), and (b) error propagation affects the later part of the generation.
This also implies that evaluating LMs solely based on short answers may not provide an adequate assessment of their factual precision, as it fails to account for errors that arise in the later stages of generation.

\myskip{
\begin{table}[t]
    \center  \mysmallfontsize %\footnotesize
    \setlength{\tabcolsep}{0.5em}
    \begin{tabular}{p{7.5cm}}
        \toprule
            \textbf{\em Single-sentence contradiction (43.3\%)} \\
            % \myblue{Gen} On November 25th, 2023, Glover Teixeira became an American citizen. \mypink{Wiki} In November 2020, Teixeira became an American citizen.\\
            \myblue{Gen} [Eric Hacker] was named the International League Pitcher of the Year. \mypink{Wiki} [Eric Hacker] was named the IL Pitcher of the Week. \\
            \myblue{Gen} She has appeared in several successful films such as (...) and Zero (2018). \mypink{Wiki}: Zero was a commercial failure. \\
        \cmidrule(lr){1-1}
            \textbf{\em Page-level contradiction (23.3\%)} \\
            \myblue{Gen} Some of [Julia Faye's] notable films include ... "Cleopatra" (1934). \mypurple{Comment} No mention of {\em Cleopatra} on the {\em Julia Faye} page, and no mention of {\em Julia Faye} on the {\em Cleopatra} page. \\
            \myblue{Gen} [Kang Ji-hwan] has donated money to various charities and organizations. \mypurple{Comment} No such mention on the {\em Kang Ji-hwan} page.
            \\
        \cmidrule(lr){1-1}
            \textbf{\em Subjective (16.7\%)} \\
            \myblue{Gen} His achievements, as an actor and as a cultural force, will surely prove to be as heroic as those of the characters he portrayed. \mypink{Wiki} Culture writer Steve Rose, in The Guardian, wrote:
            ``Chadwick Boseman began his career playing African American icons and pioneers; he ends it as one himself. His [...] achievements, as an actor and as a cultural force, will surely prove to be as heroic as those of the characters he portrayed.''\\
        \bottomrule
    \end{tabular}
    \caption{Top three categories of precision errors (\nsLabel) from PerplexityAI (Section~\ref{subsec:annotation-results}). See Appendix~\ref{subsec:annotation-qualitative} for the full results.
    \myblue{Gen} indicates the generation from PerplexityAI, and \mypink{Wiki} indicates evidence text from Wikipedia. \mypurple{Comment} indicates our comments.
    }\label{tab:ppl-error-analysis-teaser}
\end{table}
}

\vspace{-.3em}
\paragraph{Qualitative analysis of \nsLabel.}
% This section provides the full categorization of \nsLabel\ cases of PerplexityAI, summarized in Section~\ref{subsec:annotation-results}. Table~\ref{tab:ppl-error-analysis} shows examples of each category.
One of the surprising findings in our empricial analysis is that a \ours\ of PerplexityAI (71.5\%) is lower than expected despite having access to the search engine. To better understand its errors, we categorize 30 random samples whose label is \nsLabel\ (Table~\ref{tab:ppl-error-analysis}).
\vspace{-.3em}
\begin{itemize}[leftmargin=12pt]\itemsep -.2em
    \item Single-sentence contradiction: A single sentence from Wikipedia provides direct contradiction to the generation, either at a word level (numbers, dates, or entities) or beyond. %. It can be further divided into cases where the error is due to (1) simple word replacement such as numbers, dates and entities, or (2) beyond, e.g., PerplexityAI copies grandfathers of someone else (relevant to the subject) instead of the subject.
    \item Page-level contradiction: Errors found after reading the entire page, often because a fact that should have been mentioned in Wikipedia if true is missing, e.g., whether the subject appears in a particular film.
    \item Subjective: Generation is subjective, often because PerplexityAI copies subjective text from Wikipedia, e.g., directly copying a quote from a journalist without realizing it.
    \item Fact is irrelevant: Generation is irrelevant to the subject due to a search error.
    \item Wiki is inconsistent \& wrong: In the example, Wikipedia indicates that the subject won one award from the film Kick, but also includes text that they won multiple awards from Kick, which is inaccurate and cited a news article that does not support the claim.
    \item Annotation error: Annotators assign incorrect labels, typically because the information is not mentioned in the subject's Wikipedia page (likely because it is insignificant). %, or because the information is mentioned in an infobox which was not considered in our annotation process.
\end{itemize}

\noindent
We also find that, although PerplexityAI provides citations to the references, citations have little correlation with factual precision. 36.0\% and 37.6\% of supported and unsupported sentences have citations, respectively. %, indicating that unsupported sentences are in fact marginally more likely to have citations, and the overall citation ratio is low. 
Together with independent findings from \citet{liu2023evaluating}, % that citations may have low precision and recall,
this indicates that commercial LMs that incorporate search and provide citations may not be as reliable as expected.

\vspace{.2em}
More analysis is provided in Appendix~\ref{subsec:annotation-qualitative}.